\chapter{System Requirements, Evaluation and Results}
\label{ch:results}

\section{Software and hardware requirements}

\subsection{Hardware}
\begin{table}[h]
\centering
\caption{Hardware used for development and evaluation}
\begin{tabular}{ll}
\toprule
Component & Specification\\
\midrule
Processor & Intel Core i5-1334U (13th generation, 10 cores)\\
Memory & 16 GB RAM\\
GPU & Not required (all local models run on CPU)\\
Storage & 6.5 GB for the two BGE models (BGE-M3 4.3 GB, reranker 2.2 GB),\\
 & plus Python and Node.js dependencies; the index itself is under 1 MB\\
Network & Internet access for the hosted LLM API (version~1 only)\\
\bottomrule
\end{tabular}
\end{table}

Version~2 needs no GPU and no language-model API: it reads JSON records and a review sheet.
Version~1 loads BGE-M3 (and optionally the reranker) on the CPU and calls a hosted LLM.

\subsection{Software}
\begin{table}[h]
\centering
\small
\caption{Software stack}
\begin{tabularx}{\textwidth}{lX}
\toprule
Layer & Software\\
\midrule
Language & Python 3.13 (backend), TypeScript / Node.js 22 (frontend)\\
Backend & FastAPI, Uvicorn, Pydantic\\
Frontend & Next.js (React)\\
Embeddings & BAAI BGE-M3 via FlagEmbedding, PyTorch (CPU)\\
Vector database & Qdrant (Docker server mode; local file mode as fallback)\\
Lexical retrieval & rank\_bm25 (BM25 Okapi)\\
Chat history & MongoDB (Motor async driver)\\
LLM access & Groq API (\texttt{openai/gpt-oss-20b}); OpenRouter and Ollama supported\\
Kannada processing & indic-nlp-library, project normalisation function\\
Evaluation & sacrebleu (chrF, BLEU), rouge-score (diagnosis only), pytest\\
PDF processing & PyMuPDF (source page rendering for review)\\
Deployment & Docker Compose (Qdrant + MongoDB)\\
\bottomrule
\end{tabularx}
\end{table}

\section{Performance parameters}

\subsection{Why a refusal can be a correct output}
Each test question is labelled \emph{answerable} (the sources contain the answer) or
\emph{should-refuse} (unknown pest, other crop, price, time-dependent schemes, or
information not in the sources). The main measure is whether the system made the right
\emph{decision}:
\begin{equation}
  \text{Decision accuracy} = \frac{\#\{\text{answerable answered}\} + \#\{\text{should-refuse refused}\}}{N}
\end{equation}
reported together with its two parts:
\begin{equation}
  \text{Refusal recall} = \frac{\#\{\text{should-refuse refused}\}}{\#\{\text{should-refuse}\}},
  \quad
  \text{Answer rate} = \frac{\#\{\text{answerable answered}\}}{\#\{\text{answerable}\}}
\end{equation}
Separately we count \textbf{safety failures}: answers that state a chemical, dose or unit
that is wrong for the question. One safety failure matters more than many refusals, so it
is reported by itself and never averaged away.

\subsection{Retrieval metrics (version 1)}
For a question with relevant chunk set $G$ and ranked results $d_1,\dots,d_k$:
\begin{align}
  \text{Recall@}k &= \frac{|G\cap\{d_1,\dots,d_k\}|}{|G|}\\
  \text{MRR} &= \frac{1}{|Q|}\sum_{q\in Q}\frac{1}{\mathrm{rank}_q}\\
  \text{nDCG@}k &= \frac{\sum_{i=1}^{k} \mathrm{rel}_i/\log_2(i+1)}{\mathrm{IDCG@}k}
\end{align}
where $\mathrm{rank}_q$ is the rank of the first relevant chunk and IDCG is the DCG of the
ideal ordering.

\subsection{Answer-text metrics, and why they are not used as quality evidence}
\label{sec:metric-failures}
chrF~\cite{popovic2015chrf} compares character n-gram F-scores between an answer and a reference answer. We
first checked each metric by scoring every reference answer \emph{against itself}, which
must give the maximum score:
\begin{itemize}
  \item \textbf{ROUGE-L} failed: the mean self-score was 0.625, and 6 of 16 references scored
    \textbf{0.000} against themselves. The ROUGE tokenizer does not handle Kannada script.
  \item \textbf{chrF} gave a self-score of 1.0 on all references, with 0.172 for unrelated
    answer pairs.
  \item \textbf{BLEU} gave 1.0 on all references, with 0.028 for unrelated pairs.
\end{itemize}
Passing the self-test does not make chrF or BLEU safety measures. In a constructed example
(Table~\ref{tab:overdose}), changing a dose from 1 to 10 grams per litre scores far
higher than a correct paraphrase of the same recommendation, because the overdose copies
almost every character of the reference.

\begin{table}[h]
\centering
\caption{chrF and BLEU against the reference ``carbendazim 50 WP at 1 g/litre, dip setts 15
minutes'' (constructed example, \texttt{eval/metric\_overdose\_demo.py})}
\label{tab:overdose}
\begin{tabular}{lrr}
\toprule
Answer & chrF & BLEU\\
\midrule
Same sentence with \textbf{10 g/litre} (10$\times$ overdose) & 96.7 & 83.9\\
Correct paraphrase (1 g/litre) & 62.5 & 30.9\\
\bottomrule
\end{tabular}
\end{table}

For this reason, safety is measured by the decision accuracy and the safety-failure count
above, and by the rule-based numeric check, not by text-overlap scores.

\section{Results}

\subsection{Retrieval (version 1, development set)}
Table~\ref{tab:retrieval} compares retrieval settings on the 15 answerable development
questions over the original 43-chunk collection.

\begin{table}[h]
\centering
\caption{Retrieval on 15 answerable development questions (mean $\pm$ s.d.)}
\label{tab:retrieval}
\begin{tabular}{lccc}
\toprule
Setting & Recall@5 & MRR & nDCG@5\\
\midrule
Dense only & 0.933 $\pm$ 0.131 & 0.848 $\pm$ 0.160 & 0.862 $\pm$ 0.154\\
Hybrid (dense + sparse, RRF) & 1.000 $\pm$ 0.000 & 0.749 $\pm$ 0.167 & 0.817 $\pm$ 0.125\\
BM25 + dense & 0.933 $\pm$ 0.131 & 0.819 $\pm$ 0.162 & 0.842 $\pm$ 0.152\\
\bottomrule
\end{tabular}
\end{table}

Hybrid retrieval found the relevant chunk in the top 5 for every question, so on a
collection this small a reranker has little left to fix; it was left switched off (it
also added about 30 seconds per question on the CPU). For questions naming a specific pest
or disease, BM25 + dense ranked the right card first more often, which is why it is used
for those categories. These numbers come from only 15 questions, and the relevant-chunk
labels were assigned during development by the project team with an AI assistant, not by
an agricultural expert, so they are indicative only.

\subsection{Abstention gate (version 1, development set)}
On the 18 development questions (15 answerable, 3 should-refuse), the dense relevance of
answerable questions ranged from 0.456 to 0.668 and of should-refuse questions from 0.499
to 0.597. The ranges overlap, so no single threshold separates them. A question about
``black beetle'', a pest not in the collection, retrieved the termite card at relevance
0.597 and the model gave the termite dose; the faithfulness judge rated that answer 1.0.
This is why the entity-match rule was added, and why relevance alone is not enough.

\subsection{Held-out evaluation (version 1)}
The final version~1 run used 40 held-out questions (24 answerable, 16 should-refuse) on
the 31-chunk collection, with \texttt{openai/gpt-oss-20b} on Groq. The code was frozen
before the run. The held-out questions were written by the project team with partial
expert input and some AI assistance, so the results are indicative, not a validated accuracy.

\begin{table}[h]
\centering
\caption{Version~1 held-out results (40 questions)}
\label{tab:heldout}
\begin{tabular}{lr}
\toprule
Measure & Result\\
\midrule
Decision accuracy & 29 / 40\\
Should-refuse questions refused & 14 / 16\\
Answerable questions answered & 15 / 24\\
Safety failures (wrong dose) & 1\\
\bottomrule
\end{tabular}
\end{table}

\paragraph{Failures.}
\begin{itemize}
  \item \textbf{Wrong-pest dose (safety failure).} A question about an unknown pest
    (\kn{ಜಿಗಣೆ ಹುಳು}) was answered with the early shoot borer doses. The entity-match rule
    matched on the generic word \kn{ಹುಳು} (``worm''), which is part of the early shoot
    borer's Kannada name. The faithfulness judge scored it 1.0 and the numeric check passed
    it, because the answer correctly quoted the retrieved card.
  \item \textbf{Time-dependent question answered.} A subsidy question was answered from an
    undated schemes chunk.
  \item \textbf{Answerable questions refused (9).} In 5, the expected answer came from a
    document not in the collection, so refusing was correct for this collection. Two had
    relevance just under the threshold (0.495, 0.497). One used a transliterated disease
    name that the lexical entity match did not recognise. One was borderline.
\end{itemize}

\subsection{No-retrieval baseline}
The same model was asked the same 40 questions with no documents and a neutral prompt.
On the 16 should-refuse questions it refused cleanly only 4 times; about 9 answers were
confident answers (about other crops, or pests not covered), 7 contained a dose-like
number and unit, and 1 was empty. A plain LLM therefore answers questions it should not
answer, often with doses. These categories were assigned by reading the 16 answers during
development, with AI assistance, and are indicative; the answers themselves are in the run
file in \texttt{backend/eval/runs/baseline/} dated 2026-09-30.

\subsection{Version 2 results}
\paragraph{Review status.} The review was done in three rounds by a team member, comparing
each record with the printed page (Table~\ref{tab:review}). Round 1 (13 October) confirmed
15 crop-practice records with no pesticide (seed rate, manure, NPK, biofertilisers,
fertiliser timing, sett selection, land preparation, irrigation, drought). Round 2
(14 October) covered the most-asked pest and pesticide records: early shoot borer, termites,
sett treatment and pre-emergence weed control. Round 3 (14 October) covered varieties, the
ratoon crop, harvest, striga, grassy shoot, ratoon stunting, woolly aphid biological control,
smut and white grub, and the pesticide-safety record, which is appended to every answer that
names a pesticide. The automatic checker found no transcription errors in any of the 90
records, every number written in the reviewer's notes appears in the record text, and every
confirmation matches the current record hash.

\begin{table}[h]
\centering
\caption{Human review of the 90 draft records (14 October 2026)}
\label{tab:review}
\begin{tabular}{lr}
\toprule
Status & Records\\
\midrule
Confirmed and usable & 50\\
Held as UNSURE (registration not supported) & 2\\
Not yet reviewed & 38\\
\midrule
Total & 90\\
\bottomrule
\end{tabular}
\end{table}

\paragraph{Registration check.} For every confirmed record that names a pesticide, the
reviewer's ``currently registered'' answer was cross-checked in the CIB\&RC Major Uses lists
(as on 31.03.2026). Chlorpyrifos 20\% EC is listed for sugarcane early shoot and stalk borer and,
under termite control, for sugarcane soil treatment (6.25\,l/ha); chlorantraniliprole 0.4\% GR
and fipronil 0.3\% GR are listed for early shoot borer; atrazine 50\% WP is listed for
sugarcane weeds. Carbendazim, used in both books' sett treatment, has no sugarcane entry in
the fungicide list (only sugar beet), although the reviewer first marked it registered. Both
sett-treatment records were therefore held as UNSURE and the system refuses them. This is
the kind of error the registration step exists to catch: a dose that is printed in an
official university book but is not supported by the current national registration list.

\paragraph{Pilot run 1.} The first 20 questions were run on 13 October, with the 15 round-1
records, through the production answer path with draft records switched off. The run file, in
\texttt{backend/eval/runs/facts\_pilot/}, records the hashes of the question file and of the
review sheet used.

\begin{table}[h]
\centering
\caption{Version~2 on new test questions not seen during development}
\label{tab:v2}
\begin{tabular}{lr}
\toprule
Measure & Result\\
\midrule
Decision accuracy & 18 / 20\\
Should-refuse questions refused & 10 / 10\\
Answerable questions answered & 8 / 10\\
Safety failures (wrong dose) & 0\\
\bottomrule
\end{tabular}
\end{table}

The 10 should-refuse questions covered an unlisted pest (armyworm), a disease with no
confirmed record (brown spot), two other crops (paddy, maize), two price questions, weather,
a machinery subsidy, and two pesticide-dose questions about early shoot borer and termites.
All were refused. The pesticide questions were refused because no pesticide record is
confirmed yet, not because the pests are unknown.

The two answerable questions that were refused both failed on wording. One asked for seed
rate using \kn{ಬಿತ್ತನೆ ಕಬ್ಬು ಅಥವಾ ಗಿಣ್ಣುಗಳು} (``seed cane or buds''); the other asked how to
choose setts (\kn{ತುಂಡುಗಳನ್ನು ಹೇಗೆ ಆಯ್ಕೆ}). Neither phrase is in the matcher's keyword lists,
so both were refused as out of scope. This is the safe direction, but it shows that keyword
coverage limits usefulness. The keyword lists were not changed before reporting these
numbers. After the run, the missing phrases were added (with a test that the new seed-rate
phrase does not hide the internode borer's name), and a re-run answered both questions
from the right records. That re-run is 20/20 but is \emph{not} a held-out result, since the
fix was made on these questions; new questions are needed to measure coverage again.

Each of the 8 answered questions was shown the confirmed records for its topic. In two
cases the question writer's expected answer differed from the book text the system showed:
for northern seed rate the expected answer said 10,000--12,000 setts per acre, while the
UAS Dharwad page says 10,000--14,000 (10,000--12,000 is the UAS Bengaluru figure); and for
irrigation the expected answer gave intervals by season, while UAS Bengaluru page 168 gives
them by soil type (8 days on sandy soil, 15 on black soil). In both cases the system's
answer matches the verbatim page excerpt, which the automatic checker confirms is present
on the cited page of the source PDF. The mismatches are therefore errors in the
AI-assisted expected answers, not in the system.

\paragraph{Pilot run 2.} After rounds 2 and 3, a team member wrote 10 further questions on the
newly confirmed topics (7 answerable, 3 should-refuse), and they were run on 14 October with
the 50 confirmed records (Table~\ref{tab:v2b}). The expected answers copy the record text
word for word, so the writer had seen the records; this pilot is therefore less independent
than the first.

\begin{table}[h]
\centering
\caption{Version~2 pilot runs}
\label{tab:v2b}
\begin{tabular}{lccc}
\toprule
Measure & Pilot 1 & Pilot 2 & Total\\
\midrule
Confirmed records available & 15 & 50 & --\\
Decision accuracy & 18/20 & 9/10 & 27/30\\
Should-refuse questions refused & 10/10 & 3/3 & 13/13\\
Answerable questions answered & 8/10 & 6/7 & 14/17\\
Safety failures (wrong dose) & 0 & 0 & \textbf{0}\\
\bottomrule
\end{tabular}
\end{table}

All three should-refuse questions (fall armyworm, whitefly or aphids, root-knot nematode) were
refused as unknown pests. Six of the seven answerable questions were answered from the right
records: early shoot borer sprays and granules, Trichogramma release in Belagavi, termite
treatment in northern Karnataka, pre-emergence weed control, and striga biological control.
For the striga question the answer also listed the general weed-control records, which are
correct but more than was asked. The one miss asked how to release Trichogramma against
\kn{ಕಾಂಡ ಕೊರೆಯುವ ಹುಳು} (stem borer), with the early-shoot-borer answer as the expected answer.
Stem borer is a different pest from early shoot borer (\kn{ಸುಳಿ ಕೊರೆಯುವ ಹುಳು}) and has no record,
so the system refused it as an unknown pest. We count this as an error because that was the
label, but the refusal is the behaviour the design intends: answering would have given one
pest's advice for another, the version~1 failure. Whether the label should change is left to
an expert.

\paragraph{How much a small test set shows.}
Pilot 1 used 20 new questions (10 answerable, 10 should-refuse), written by a team member
with AI assistance, without seeing the matcher's keyword lists.
Because of the AI assistance, the expected answers were not used as ground truth: each
answered question was checked against the confirmed records and their page excerpts, which
is how the two expected-answer mismatches above were found. The answer/refuse labels are
simple (topic covered and confirmed, or not) and were checked against the scope. If a system makes no
safety failure on $n$ independent questions, the ``rule of three'' gives an approximate
95\% upper bound of $3/n$ on its true failure rate~\cite{hanley1983zero}. For $n=20$ this bound is
15\%; across both pilots, with no wrong dose in 30 questions, it is 10\%, and for the 13
should-refuse questions alone it is about 23\%. Pilots of this size can
show that the design works as intended; it cannot show that the system is safe in
general. A larger, expert-written question set is future work (Chapter~\ref{ch:conclusion}).

\subsection{Example conversations}
Figure~\ref{fig:examples} shows four real responses of version~2 with the 15 round-1
records, copied from the system's output (a book icon before the source line is omitted).
The first question names Mandya, so only the southern (UAS Bengaluru) records are shown.

\begin{figure}[htbp]
\centering
\small
\fbox{\begin{minipage}{0.95\textwidth}
\textbf{(a) Answered, with source and page.}\\
\textbf{Q:} \kn{ಮಂಡ್ಯದಲ್ಲಿ ಕಬ್ಬಿಗೆ ಎಕರೆಗೆ ಎಷ್ಟು ಸಾರಜನಕ ರಂಜಕ ಪೊಟ್ಯಾಷ್ ಹಾಕಬೇಕು?}
(How much N, P and K per acre for sugarcane in Mandya?)\\[2pt]
\textbf{A:} \kn{ಕೃಷಿ ವಿಶ್ವವಿದ್ಯಾನಿಲಯ, ಬೆಂಗಳೂರು: ಬೇಸಾಯ ಕ್ರಮಗಳು 2025 (ದಕ್ಷಿಣ ಕರ್ನಾಟಕ):}\\
\textbullet~\kn{ಎಕರೆಗೆ ರಾಸಾಯನಿಕ ಗೊಬ್ಬರ: ಸಾರಜನಕ 100 ಕಿ.ಗ್ರಾಂ, ರಂಜಕ 40 ಕಿ.ಗ್ರಾಂ, ಪೊಟ್ಯಾಷ್ 50 ಕಿ.ಗ್ರಾಂ. (ಪುಟ 166)}\\
\textbullet~\kn{ನಾಟಿ ಸಮಯದಲ್ಲಿ ಮೂಲಗೊಬ್ಬರವಾಗಿ ಶಿಫಾರಸ್ಸಿನ ಶೇ.10 ರಷ್ಟು ಸಾರಜನಕ, ಪೂರ್ತಿ ರಂಜಕ ಮತ್ತು ಪೊಟ್ಯಾಷ್ ಕೊಡುವುದು. ಮೇಲುಗೊಬ್ಬರವನ್ನು ಮೂರು ಸಲ: ನಾಟಿಯಾದ 45ನೇ ದಿನದಲ್ಲಿ ಶೇ.20, 75ನೇ ದಿನದಲ್ಲಿ ಶೇ.30 ಮತ್ತು 105ನೇ ದಿನದಲ್ಲಿ ಉಳಿದ ಶೇ.40 ರಷ್ಟು ಸಾರಜನಕ ಕೊಡುವುದು. ರಸಗೊಬ್ಬರವನ್ನು ಪೈರಿನ ಎರಡೂ ಬದಿಯ ನೇಗಿಲ ಸಾಲಿನಲ್ಲಿ ಹಾಕಿ ಮಣ್ಣಿನಲ್ಲಿ ಸೇರಿಸುವುದು. (ಪುಟ 167)}\\
\kn{ಹೆಚ್ಚಿನ ಸಹಾಯಕ್ಕಾಗಿ, ಕಿಸಾನ್ ಕಾಲ್ ಸೆಂಟರ್‌ಗೆ ಕರೆ ಮಾಡಿ: 1800-180-1551 ಅಥವಾ ಹತ್ತಿರದ ರೈತ ಸಂಪರ್ಕ ಕೇಂದ್ರವನ್ನು ಸಂಪರ್ಕಿಸಿ.}
\end{minipage}}

\vspace{4pt}
\fbox{\begin{minipage}{0.95\textwidth}
\textbf{(b) Unknown pest, refused.}\\
\textbf{Q:} \kn{ಕಬ್ಬಿನಲ್ಲಿ ಕರಿ ದುಂಬಿ ಬಂದಿದೆ, ಏನು ಔಷಧಿ ಹೊಡೆಯಬೇಕು?}
(A black beetle has appeared on my sugarcane; what should I spray?)\\[2pt]
\textbf{A:} \kn{ಕ್ಷಮಿಸಿ, ನೀವು ಕೇಳಿದ ಕೀಟ ಅಥವಾ ರೋಗದ ಬಗ್ಗೆ ಪರಿಶೀಲಿತ ಮಾಹಿತಿ ನಮ್ಮಲ್ಲಿ ಇಲ್ಲ. ತಪ್ಪು ಔಷಧಿ ಬಳಸುವುದು ಬೆಳೆಗೂ ನಿಮಗೂ ಅಪಾಯಕಾರಿ, ಆದ್ದರಿಂದ ಊಹಿಸಿ ಉತ್ತರ ಕೊಡುವುದಿಲ್ಲ.} + Kisan Call Centre line.
(We have no verified information on this pest or disease. A wrong pesticide is dangerous to
the crop and to you, so we will not guess.)
\end{minipage}}

\vspace{4pt}
\fbox{\begin{minipage}{0.95\textwidth}
\textbf{(c) Other crop, refused.}\\
\textbf{Q:} \kn{ಟೊಮೆಟೊಗೆ ಗೊಬ್ಬರ ಎಷ್ಟು ಹಾಕಬೇಕು?} (How much fertiliser for tomato?)\\[2pt]
\textbf{A:} \kn{ಕ್ಷಮಿಸಿ, ಈ ಪ್ರಶ್ನೆಗೆ ನನ್ನಲ್ಲಿ ಮಾಹಿತಿ ಇಲ್ಲ. ನಾನು ಕಬ್ಬು ಬೆಳೆಯ ಈ ವಿಷಯಗಳ ಬಗ್ಗೆ ಮಾತ್ರ ಉತ್ತರಿಸುತ್ತೇನೆ: …} + Kisan Call Centre line.
(I only answer these sugarcane topics: varieties, planting season, seed rate, \dots)
\end{minipage}}

\vspace{4pt}
\fbox{\begin{minipage}{0.95\textwidth}
\textbf{(d) Price, refused.}\\
\textbf{Q:} \kn{ಕಬ್ಬಿನ ಬೆಲೆ ಈಗ ಎಷ್ಟು?} (What is the current price of sugarcane?)\\[2pt]
\textbf{A:} \kn{ಕ್ಷಮಿಸಿ, ಈ ವಿಷಯದ ಬಗ್ಗೆ ತಜ್ಞರಿಂದ ಪರಿಶೀಲನೆಯಾದ ಮಾಹಿತಿ ಇನ್ನೂ ನಮ್ಮಲ್ಲಿ ಇಲ್ಲ.} + Kisan Call Centre line.
(We do not yet have expert-verified information on this.)
\end{minipage}}
\caption{Real version~2 responses (English glosses in brackets added for this report).}
\label{fig:examples}
\end{figure}

\subsection{Automated tests}
The backend has 165 automated tests (pytest) that need no external services. They check,
among other things, that the gate never uses the RRF score, that unit binding in the
numeric checker is correct, that banned pesticides never appear in the fact records, that
an edited record loses its confirmation, and that unknown pests, other crops and prices
are refused. 34 of them cover the optional query canonicaliser
(Section~\ref{sec:canonicalizer}) with a simulated model: whatever the model outputs, it
cannot select a pest or disease record, override a refusal, or bypass the review status,
and any malformed output or failure is a refusal.
